\documentclass[12pt,a4paper]{article}
\usepackage{amsmath,amssymb,amsthm}
\usepackage{geometry}
\usepackage{enumitem}
\usepackage{xcolor}
\usepackage{mdframed}
\usepackage{multicol}
\usepackage{array}
\geometry{margin=2cm}

\definecolor{anscolor}{RGB}{0,120,0}
\definecolor{partbg}{RGB}{230,240,255}
\newcommand{\ans}[1]{\quad\textcolor{anscolor}{\textbf{[Answer: #1]}}}
\newcommand{\shead}[1]{\vspace{1em}{\large\bfseries\color{blue}#1}\vspace{0.3em}\hrule\vspace{0.5em}}

\begin{document}

\begin{center}
{\LARGE\bfseries CSIR-UGC NET June 2025}\\[4pt]
{\large Mathematical Sciences (Code 704) \quad 28 July 2025, Shift 2}\\[4pt]
{\large Complete Question Paper with Final Answer Key}
\end{center}
\hrule\vspace{1em}

%% ============================================================
\shead{PART A --- General Aptitude (Q.1--Q.20)}
%% ============================================================
\begin{mdframed}[backgroundcolor=partbg,linewidth=0pt]
Marking: $+2$ correct, $-0.5$ wrong. Attempt any 15 of 20.
\end{mdframed}
\vspace{0.5em}

\textbf{Q.1} Consider the following statements:\\
Statement I: All Booklets are Manuals.\\
Statement II: All Manuals are Catalogues.\\
If Statements I and II are True, which one of the following conclusions can be conclusively drawn?
\begin{enumerate}[label=\arabic*.]
\item All Manuals are Booklets. \item All Catalogues are Booklets.
\item All Booklets are Catalogues. \item All Catalogues are Manuals.
\end{enumerate}\ans{3}

\textbf{Q.2} The given Venn diagram shows numbers of players playing one or more than one sport (Football, Cricket, Tennis). The percentage of players who play exactly two sports is closest to \underline{\hspace{1cm}}\%.
\begin{enumerate}[label=\arabic*.]
\item 5 \item 14 \item 28 \item 32
\end{enumerate}\ans{3}

\textbf{Q.3} The value of a company is measured as the total value of its shares owned by different investors. Rakesh owns 2/15 of the shares of a company. He sells 1/3 of his shares for Rs.\,75,000/-. What is the total value of the company at that time?
\begin{enumerate}[label=\arabic*.]
\item Rs.\,15,75,800 \item Rs.\,16,87,500 \item Rs.\,17,75,800 \item Rs.\,18,27,500
\end{enumerate}\ans{2}

\textbf{Q.4} A car has wheels of diameter 36 cm. If it runs at a speed of 60 km/h, then the rotation per minute (RPM) will be closest to \underline{\hspace{1cm}}.
\begin{enumerate}[label=\arabic*.]
\item 884 \item 898 \item 906 \item 986
\end{enumerate}\ans{1}

\textbf{Q.5} A cylindrical container of radius 20 cm was filled with water up to 25 cm height. A solid spherical ball of radius 7 cm was then immersed in the water. What would be the approximate increase in water level in the container after the ball was fully immersed?
\begin{enumerate}[label=\arabic*.]
\item 1.14 cm \item 2.28 cm \item 5.50 cm \item 7.00 cm
\end{enumerate}\ans{1}

\textbf{Q.6} The market share (\%) and annual production of scooters from five automobile companies P, Q, R, S, and T are shown in graphs. If the profit of a company is directly proportional to the ratio of market share to production, then which of the following statements is/are CORRECT?\\
Statement X: Companies T and P have same profit\\
Statement Y: Company R has the maximum profit\\
Statement Z: Company S has the minimum profit
\begin{enumerate}[label=\arabic*.]
\item X and Y \item X and Z \item Y and Z \item Only Z
\end{enumerate}\ans{3}

\textbf{Q.7} Rahul and his father started jogging on a circular track of radius $r$ ($r>2$). Rahul completed one round and stopped. His father got tired half way into the first round and returned to his starting point along a straight line. What is the ratio of the distances covered by Rahul and his father?
\begin{enumerate}[label=\arabic*.]
\item $\pi r/(\pi+2)$ \item $2\pi/(\pi+2)$ \item $1$ \item $2$
\end{enumerate}\ans{2}

\textbf{Q.8} Kavita starts from her house and walks 200 m northward, then turns $45°$ right and walks 70 m. After that, she turns $90°$ right and walks 70 m. Which of the following is the closest value of the shortest distance between Kavita's current location and her house?
\begin{enumerate}[label=\arabic*.]
\item 296 m \item 240 m \item 200 m \item 223 m
\end{enumerate}\ans{4}

\textbf{Q.9} The initial monthly salaries of employees John, Riya, and Sunil were in the proportion 4:3:5. After an increase of Rs.\,10000 monthly to all, the new proportion becomes 6:5:7. What was the initial salary of Sunil?
\begin{enumerate}[label=\arabic*.]
\item Rs.\,20000 \item Rs.\,25000 \item Rs.\,30000 \item Rs.\,35000
\end{enumerate}\ans{2}

\textbf{Q.10} Numbers of Rose, Lotus, and Marigold plants in a garden are in the proportion 8:5:7. Later, 75\%, 40\% and 50\% more plants of their respective categories were added. What will be the new proportion of plants, in the same order?
\begin{enumerate}[label=\arabic*.]
\item 5:3:4 \item 4:2:3 \item 5:4:3 \item 7:4:5
\end{enumerate}\ans{2}

\textbf{Q.11} What will be the digit at the unit's place of $1^3+2^3+3^3+4^3+5^3+6^3+7^3+8^3+9^3$?
\begin{enumerate}[label=\arabic*.]
\item 0 \item 5 \item 7 \item 9
\end{enumerate}\ans{2}

\textbf{Q.12} Suresh asked Ramesh to identify the person in a photo that the latter is holding. Ramesh responds, ``I have no brothers or sisters. However, that man's father is my father's son.'' Who is the person in the photo?
\begin{enumerate}[label=\arabic*.]
\item Suresh \item Ramesh \item Ramesh's son \item Ramesh's cousin
\end{enumerate}\ans{3}

\textbf{Q.13} Three identical semi-circles are arranged as shown. (The total horizontal span is $15+10+15=40$ and the overlap is 10.) What is the diameter of the semi-circles?
\begin{enumerate}[label=\arabic*.]
\item $5\pi$ \item $20$ \item $15\pi/2$ \item $25$
\end{enumerate}\ans{2}

\textbf{Q.14} A number is mistakenly divided by 2 instead of being multiplied by 2. What is the change in the result caused by this mistake?
\begin{enumerate}[label=\arabic*.]
\item 25\% \item 50\% \item 75\% \item 100\%
\end{enumerate}\ans{3}

\textbf{Q.15} The following diagram represents the relationship between four categories. The categories could be
\begin{enumerate}[label=\arabic*.]
\item Rivers, water bodies, oceans, sources of evaporation
\item Parliamentarians, celebrities, elected persons, professional politicians
\item Monkeys, four-legged animals, pet animals, land animals
\item Furniture, chairs, seats, wooden objects
\end{enumerate}\ans{4}

\textbf{Q.16} In a code, the word DELTOID is written as 3152893. Then LOTION could be written as
\begin{enumerate}[label=\arabic*.]
\item 582986 \item 582981 \item 198396 \item 198392
\end{enumerate}\ans{1}

\textbf{Q.17} Sum of the digits of a two-digit number `ab' is subtracted from the number and the result is divided by 9. Then the result of this will be
\begin{enumerate}[label=\arabic*.]
\item always a \item always b \item neither a nor b \item either a or b depending on a+b
\end{enumerate}\ans{1}

\textbf{Q.18} A circle of radius 1 unit is divided into four quarters and rejoined as shown below. What is the area of this shape?
\begin{enumerate}[label=\arabic*.]
\item $\pi$ \item $1$ \item $2$ \item $4$
\end{enumerate}\ans{3}

\textbf{Q.19} A stock market trader has lost two thirds of her investment on a day. Next day she recovered one third of the previous day's loss. What fraction of her initial investment is she left with?
\begin{enumerate}[label=\arabic*.]
\item $\dfrac{1}{3}$ \item $\dfrac{2}{3}$ \item $\dfrac{2}{9}$ \item $\dfrac{5}{9}$
\end{enumerate}\ans{4}

\textbf{Q.20} Three friends, Mr.\ Rahman, Mr.\ George and Mr.\ Vedant, met after a long time. They were wearing red, green and violet colour shirts. Mr.\ Rahman and the person wearing violet shirt noticed that none of the three is wearing a colour that starts with same letter as his name. Which one of the following is the correct match of the persons with the colour of their shirts?
\begin{enumerate}[label=\arabic*.]
\item Rahman-Violet, George-Red, Vedant-Green
\item Rahman-Green, George-Violet, Vedant-Red
\item Rahman-Green, George-Red, Vedant-Violet
\item Rahman-Red, George-Violet, Vedant-Green
\end{enumerate}\ans{2}

%% ============================================================
\shead{PART B --- Mathematics MCQ (Q.21--Q.60)}
%% ============================================================
\begin{mdframed}[backgroundcolor=partbg,linewidth=0pt]
Marking: $+3$ correct, $-0.75$ wrong. Attempt any 25 of 40.
\end{mdframed}
\vspace{0.5em}

\textbf{Q.21} Let $p,q$ be non-negative integers. Consider the following statements:\\
(A) There is an integer $k\ge1$ such that $p+k=q$.\\
(B) There is an integer $k\ge1$ such that $q+k=p$.\\
Which of the following statements is true?
\begin{enumerate}[label=\arabic*.]
\item There exist non-negative integers $p,q$ such that both (A) and (B) are true.
\item Both (A) and (B) are false if and only if $p=q$.
\item For all non-negative integers $p$ and $q$, (A) or (B) is true.
\item There exists $p\neq q$ such that both (A) and (B) are false.
\end{enumerate}\ans{2}

\textbf{Q.22} Let
\begin{align*}
A&=\Bigl\{\tfrac{p}{q}\in(0,1):p\in\mathbb{N},\,q=2^n\text{ for some }n\in\mathbb{N}\cup\{0\},\,\gcd(p,q)=1\Bigr\},\\
B&=\Bigl\{\tfrac{p}{q}\in(0,1):p\in\mathbb{N},\,q=2^n5^m\text{ for some }n,m\in\mathbb{N}\cup\{0\},\,\gcd(p,q)=1\Bigr\},\\
C&=\Bigl\{\tfrac{p}{q}\in(0,1):\tfrac{p}{q}\text{ has terminating decimal expansion}\Bigr\}
\end{align*}
be subsets of $(0,1)$. Which of the following is true?
\begin{enumerate}[label=\arabic*.]
\item $A\subsetneq C$ and $B\not\subsetneq C$
\item $A\subsetneq C\subsetneq B$
\item $A\subsetneq B\not\subsetneq C$
\item $A\subsetneq B=C$
\end{enumerate}\ans{4}

\textbf{Q.23} Let $A,B$ be non-empty subsets of $\mathbb{N}$ with cardinality $|A|\ge2$. Let
\[S_1=\{f:A\to B\mid f\text{ is one-to-one}\}\quad\text{and}\quad S_2=\{g:B\to A\mid g\text{ is onto}\}.\]
Which of the following is true?
\begin{enumerate}[label=\arabic*.]
\item If $A\subsetneq B$ and $B$ is finite, then there is a one-to-one map from $S_2$ to $S_1$.
\item If $B=\mathbb{N}$, then there exists a one-to-one map from $S_2$ to $B$.
\item If $B=\mathbb{N}$ and $A$ is finite, then there exists a one-to-one map from $B$ to $S_1$.
\item If $A$ is finite, then $S_2$ is finite for any $B$.
\end{enumerate}\ans{3}

\textbf{Q.24} Let $f:\mathbb{R}\setminus\mathbb{Q}\to\mathbb{R}\setminus\mathbb{Q}$ be the function defined as $f(x)=\dfrac{3x+2}{4x+3}$. Let $x_1\in\mathbb{R}\setminus\mathbb{Q}$. For $n\ge1$, define $x_{n+1}=f(x_n)$. Suppose that the sequence $(x_n)_{n\ge1}$ converges to a real number $\ell$. Which of the following is true?
\begin{enumerate}[label=\arabic*.]
\item If $\ell$ is positive, then $\ell=\dfrac{\sqrt{3}}{2}$.
\item If $\ell$ is positive, then $\ell=\dfrac{1}{\sqrt{2}}$.
\item If $\ell$ is negative, then $\ell=-\dfrac{\sqrt{3}}{2}$.
\item If $\ell$ is negative, then $\ell=-\dfrac{1}{2}$.
\end{enumerate}\ans{2}

\textbf{Q.25} Let $f(x)=x\log_e\!\left(1+\dfrac{1}{x}\right)$ for $x\in(0,\infty)$. Which of the following is true?
\begin{enumerate}[label=\arabic*.]
\item $f$ is unbounded. \item $f$ is increasing.
\item $\lim_{x\to\infty}f(x)=2$. \item $f$ is decreasing.
\end{enumerate}\ans{2}

\textbf{Q.26} For each $n\ge1$, let $f_n:[0,1]\to\mathbb{R}$ be defined as
\[f_n(x)=\begin{cases}nx&\text{if }x\in\bigl[0,\tfrac{1}{n}\bigr],\\2-nx&\text{if }x\in\bigl(\tfrac{1}{n},\tfrac{2}{n}\bigr],\\0&\text{if }x\in\bigl(\tfrac{2}{n},1\bigr].\end{cases}\]
Which of the following is true?
\begin{enumerate}[label=\arabic*.]
\item $(f_n)_{n\ge1}$ converges uniformly on $[0,1]$ to a continuous function $f$.
\item $(f_n)_{n\ge1}$ converges pointwise on $[0,1]$ to a discontinuous function $f$.
\item $(f_n)_{n\ge1}$ converges pointwise on $[0,1]$ to a continuous function $f$.
\item $(f_n)_{n\ge1}$ does not converge pointwise on $[0,1]$.
\end{enumerate}\ans{3}

\textbf{Q.27} Which of the following polynomials is the characteristic polynomial of a real $2\times2$ matrix $A$ such that $\text{trace}(A)=7$ and $\text{trace}(A^2)=29$?
\begin{enumerate}[label=\arabic*.]
\item $t^2+7t+10$ \item $t^2-7t+29$ \item $t^2-7t-10$ \item $t^2-7t+10$
\end{enumerate}\ans{4}

\textbf{Q.28} Let $\mathbb{F}_5$ denote the field with 5 elements. How many $2\times2$ matrices with entries in $\mathbb{F}_5$ have rank one?
\begin{enumerate}[label=\arabic*.]
\item 125 \item 144 \item 145 \item 480
\end{enumerate}\ans{2}

\textbf{Q.29} Let $X$ be the $\mathbb{R}$-vector space of all twice differentiable real valued functions on $[0,1]$. Consider the linear map $\phi:X\to\mathbb{R}^3$ defined by $\phi(f)=(f(1),f'(1),f''(1))$. Which of the following is true?
\begin{enumerate}[label=\arabic*.]
\item The dimension of $X/\ker\phi$ is 3.
\item $\ker\phi$ is finite dimensional.
\item The dimension of $X/\ker\phi$ is 1.
\item $X$ is finite dimensional.
\end{enumerate}\ans{1}

\textbf{Q.30} Consider the real matrix $A=\begin{pmatrix}29&0&55&17\\1&28&46&26\\17&13&33&38\\21&67&0&13\end{pmatrix}$. What is the largest real eigenvalue of $A$?
\begin{enumerate}[label=\arabic*.]
\item 101 \item 67 \item 103 \item 113
\end{enumerate}\ans{1}

\textbf{Q.31} Let $C[0,\pi]$ be the real vector space of real-valued continuous functions on the closed interval $[0,\pi]$. For positive integers $n$, define $f_n\in C[0,\pi]$ by
\[f_n(x)=\begin{cases}\dfrac{\sin(nx)}{\sin x}&\text{if }x\in(0,\pi),\\ n&\text{if }x=0,\\ (-1)^{n-1}n&\text{if }x=\pi.\end{cases}\]
Let $V$ be the real subspace of $C[0,\pi]$ spanned by $\{f_1,f_2,f_3\}$. Consider the inner product on $V$ given by $\langle f,g\rangle=\dfrac{2}{\pi}\int_0^\pi f(x)g(x)\sin^2x\,dx$. Which of the following is true?
\begin{enumerate}[label=\arabic*.]
\item $f_4\in V$
\item $\{f_1,f_2,f_3\}$ is an orthonormal basis of $V$.
\item The dimension of $V$ is 2.
\item $\{f_1,f_2,f_3\}$ is an orthogonal set but not orthonormal.
\end{enumerate}\ans{2}

\textbf{Q.32} Let $V=\{ax^3+bx^2+cx\mid a,b,c\in\mathbb{R}\}$. For $f\in V$, define $Q(f)=\int_{-1}^1(f'(t))^2\,dt$, where $f'$ denotes the derivative of $f$. Which of the following is \textbf{FALSE}?
\begin{enumerate}[label=\arabic*.]
\item $Q$ is a positive definite quadratic form on $V$.
\item $Q$ takes every positive real value.
\item $Q(x)=2$.
\item For all $f,g\in V$, $Q(f+g)=Q(f)+Q(g)$.
\end{enumerate}\ans{4}

\textbf{Q.33} Let $f:\mathbb{C}\to\mathbb{C}$ be a polynomial map. For $R>0$, let $\gamma_R:[0,1]\to\mathbb{C}$ be the map $t\mapsto Re^{2\pi it}$. Suppose that there exists $c\in\mathbb{R}$ such that $\int_0^1|(f\circ\gamma_R)(t)\gamma_R'(t)|\,dt\to c$ as $R\to\infty$. Which of the following is \textbf{FALSE}?
\begin{enumerate}[label=\arabic*.]
\item The function $zf(1/z)\to0$ as $|z|\to\infty$.
\item The function $f$ is constant.
\item $c=0$.
\item $c>0$.
\end{enumerate}\ans{4}

\textbf{Q.34} Let $X$ be the image of the interval $[0,1]$ under the M\"{o}bius transformation $f(z)=\dfrac{z-i}{z+i}$. Which of the following is true?
\begin{enumerate}[label=\arabic*.]
\item $X$ is the line segment joining $-1$ and $-i$.
\item $X=\{e^{i\theta}\mid\theta\in[\pi,\tfrac{3\pi}{2}]\}$.
\item $X$ is the line segment joining $-1$ to $1$.
\item $X=\{e^{i\theta}\mid\theta\in[-\tfrac{\pi}{2},\pi]\}$.
\end{enumerate}\ans{2}

\textbf{Q.35} Which of the following is true?
\begin{enumerate}[label=\arabic*.]
\item There exists an entire function $f$ such that $f^{(n)}(0)=\dfrac{n!}{n^n}$ for all positive integers $n$.
\item There exists an entire function $f$ such that $f^{(n)}(0)=n!\,n^n$ for all positive integers $n$.
\item There exists an entire function $f$ such that $f^{(n)}(0)=(n-1)!$ for all positive integers $n$.
\item There exists an entire function $f$ such that $f^{(n)}(0)=n!\,n$ for all positive integers $n$.
\end{enumerate}\ans{1}

\textbf{Q.36} Let $f$ be an entire function such that $f(\mathbb{C})\subset\{x+iy\mid y=x+1\}$. Which of the following is true?
\begin{enumerate}[label=\arabic*.]
\item $|f(z)|\to\infty$ as $|z|\to\infty$. \item $\dfrac{f(z)}{z}\to0$ as $|z|\to\infty$.
\item $zf(z)\to0$ as $|z|\to\infty$. \item $f(z)\to0$ as $|z|\to\infty$.
\end{enumerate}\ans{2}

\textbf{Q.37} Which of the following is true?
\begin{enumerate}[label=\arabic*.]
\item $p\nmid1+(p-1)!$ for some odd prime $p$.
\item $p\mid(1234)^{p-1}-1$ for all primes $p>700$.
\item There exist $a\in\mathbb{Z}$ and a prime $p>11$ such that $p\nmid a^p-a$.
\item $p\nmid\dfrac{(p^2)!}{(p!)^2}$ for some odd prime $p$.
\end{enumerate}\ans{2}

\textbf{Q.38} Let $A$ be a subring of the field of rationals $\mathbb{Q}$ such that for any nonzero rational $r\in\mathbb{Q}$, $r\in A$ or $1/r\in A$. Which of the following is \textbf{FALSE}?
\begin{enumerate}[label=\arabic*.]
\item The set $\bigl\{a\in A:\tfrac{1}{a}\notin A\bigr\}\cup\{0\}$ is an additive subgroup of $\mathbb{Q}$.
\item $A$ has at most one maximal ideal.
\item If $A\neq\mathbb{Q}$, then $A$ has infinitely many prime ideals.
\item For any nonzero $a,b\in A$, $a$ divides $b$ or $b$ divides $a$ in $A$.
\end{enumerate}\ans{3}

\textbf{Q.39} Which of the following is true?
\begin{enumerate}[label=\arabic*.]
\item The ideal $2\mathbb{Z}[i]$ is maximal in $\mathbb{Z}[i]$.
\item The ideal $X\mathbb{C}[X,Y]$ is maximal in $\mathbb{C}[X,Y]$.
\item The set of all polynomials in $\mathbb{C}[X]$ whose coefficients add up to 0 is a maximal ideal in $\mathbb{C}[X]$.
\item The ideal $(\sqrt{2}-1)\mathbb{Z}[\sqrt{2}]$ is maximal in $\mathbb{Z}[\sqrt{2}]$.
\end{enumerate}\ans{3}

\textbf{Q.40} Let $S=\{1,2,3,4,5\}$ be equipped with the topology $\tau=\{\emptyset,\{1\},S\}$. What is the number of homeomorphisms of $S$ onto itself?
\begin{enumerate}[label=\arabic*.]
\item 25 \item 120 \item 24 \item 6
\end{enumerate}\ans{3}

\textbf{Q.41} If $\varphi(x)=x$ is a solution of the ordinary differential equation (ODE)
\[\dfrac{d^2y}{dx^2}-\left(\dfrac{2}{x^2}+\dfrac{1}{x}\right)\!\left(x\dfrac{dy}{dx}-y\right)=0,\quad0<x<\infty,\]
then the general solution of the ODE is given by
\begin{enumerate}[label=\arabic*.]
\item $(a+be^{-2x})x,\;a,b\in\mathbb{R}$
\item $(a+be^{2x})x,\;a,b\in\mathbb{R}$
\item $ae^x+bx,\;a,b\in\mathbb{R}$
\item $(a+be^x)x,\;a,b\in\mathbb{R}$
\end{enumerate}\ans{4}

\textbf{Q.42} Let $(\lambda_n)_{n\in\mathbb{N}}$ be the sequence of eigenvalues of the Sturm--Liouville problem
\[\dfrac{d}{dx}\!\left(x\dfrac{dy}{dx}\right)+\dfrac{\lambda}{x}y=0,\quad1<x<e^{2\pi},\quad y(1)=0,\;y(e^{2\pi})=0.\]
Then $\displaystyle\sum_{n=1}^\infty\dfrac{1}{\lambda_n}$ is equal to
\begin{enumerate}[label=\arabic*.]
\item $\dfrac{\pi^2}{12}$ \item $\dfrac{2\pi^2}{3}$ \item $\dfrac{\pi^2}{4}$ \item $\dfrac{\pi^2}{16}$
\end{enumerate}\ans{2}

\textbf{Q.43} Let $u=u(x,y)$ be the solution to the Cauchy problem
\[(y+u)\dfrac{\partial u}{\partial x}+y\dfrac{\partial u}{\partial y}=x-y,\quad x\in\mathbb{R},\;y>0,\quad u(x,1)=1+x,\quad x\in\mathbb{R}.\]
Then which of the following is true?
\begin{enumerate}[label=\arabic*.]
\item $u(1,1)=2$ \item $u(2,2)=4$
\item $u(3,3)=\dfrac{3}{2}$ \item $u(4,4)=\dfrac{2}{3}$
\end{enumerate}\ans{1}

\textbf{Q.44} Let $u=u(x,t)$ be the solution of
\[\dfrac{\partial^2u}{\partial t^2}-\dfrac{\partial^2u}{\partial x^2}=0,\;x\in\mathbb{R},\;t>0;\quad u(x,0)=1+x^2,\;x\in\mathbb{R};\quad\dfrac{\partial u}{\partial t}(x,0)=x+1,\;x\in\mathbb{R}.\]
Then the value of $u(1,1)$ is
\begin{enumerate}[label=\arabic*.]
\item 2 \item 3 \item 4 \item 5
\end{enumerate}\ans{4}

\textbf{Q.45} If the function $s:[0,4]\to\mathbb{R}$ defined by
\[s(x)=\begin{cases}a(x-2)^2+b(x-1)^2, & 0\le x\le1,\\(x-2)^2, & 1<x\le3,\\2c(x-2)^2+(x-3)^3, & 3<x\le4\end{cases}\]
is a cubic spline, then the value of $2a+b+2c$ is
\begin{enumerate}[label=\arabic*.]
\item 2 \item 3 \item 4 \item 5
\end{enumerate}\ans{2}

\textbf{Q.46} Let $y(x)$ be the extremal of the functional
\[J[y]=\int_0^{\pi/4}\bigl((y')^2-4y^2+2xy\bigr)\,dx\]
subject to $y(0)=0$, $y\!\left(\dfrac{\pi}{4}\right)=1$. Then $y(x)$ is equal to
\begin{enumerate}[label=\arabic*.]
\item $\left(1-\dfrac{\pi}{4}\right)\sin(2x)+x$
\item $\left(1-\dfrac{\pi}{16}\right)\sin(2x)+\dfrac{x}{4}$
\item $\left(1+\dfrac{\pi}{4}\right)\sin(2x)-x$
\item $\left(1+\dfrac{\pi}{16}\right)\sin(2x)-\dfrac{x}{4}$
\end{enumerate}\ans{2}

\textbf{Q.47} If $y(x)$ is the solution of the integral equation
\[y(x)=x^2+2\int_0^1xty(t)\,dt,\]
then which of the following is true?
\begin{enumerate}[label=\arabic*.]
\item $y(0)+y(1)=\dfrac{1}{2}$
\item $y(-1)+y(1)=1$
\item $y'(0)+y'(1)=\dfrac{3}{2}$
\item $y'(-1)+y'(1)=3$
\end{enumerate}\ans{4}

\textbf{Q.48} Suppose a dynamical system has the Lagrangian $L=(\dot{q}_1)^2+(\dot{q}_2)^2+(q_1)^2+\dot{q}_1\dot{q}_2$. If $p_1$ and $p_2$ are momenta conjugate to $q_1$ and $q_2$, respectively, then which of the following is true?
\begin{enumerate}[label=\arabic*.]
\item $\dot{p}_1=2q_1$, $\dot{p}_2=0$
\item $\dot{p}_1=-q_1$, $\dot{p}_2=0$
\item $\dot{p}_1=-\dfrac{q_1}{2}$, $p_2=q_2$
\item $\dot{p}_1=q_1$, $p_2=-q_2$
\end{enumerate}\ans{1}

\textbf{Q.49} Suppose that we have a data set consisting of $2n+1$ observations for some $n\in\mathbb{N}$. Value of each observation is either $x$ or $x+r$, where $x\in\mathbb{N}$, $r\ge0$. Then, which of the following statements is always true?
\begin{enumerate}[label=\arabic*.]
\item The mean and median of the data will be different if and only if $r>0$
\item Variance of the data is positive if and only if $r>0$
\item Mean and mode of the data will be same if and only if $r=0$
\item Median and mode of the data will be same for all values of $r\ge0$
\end{enumerate}\ans{4}

\textbf{Q.50} A biased six-faced die is tossed once. Suppose that the probability of any prime number showing up is twice that of any non-prime number showing up. Then, the probability that an odd number will show up is
\begin{enumerate}[label=\arabic*.]
\item $\dfrac{1}{3}$ \item $\dfrac{2}{3}$ \item $\dfrac{4}{9}$ \item $\dfrac{5}{9}$
\end{enumerate}\ans{4}

\textbf{Q.51} Let $Z_1,Z_2,\ldots$ be a sequence of independent and identically distributed random variables having discrete uniform distribution over $\{1,2,\ldots,2024\}$. Let $Y_n=\sum_{i=1}^nZ_i$, $n\ge2$. Further, let $X_n$ be the remainder when $Y_n$ is divided by 2025. Then, which of the following is true?
\begin{enumerate}[label=\arabic*.]
\item $\lim_{n\to\infty}P(X_n=0)=\dfrac{1}{2026}$
\item $\lim_{n\to\infty}P(X_n=0)=\dfrac{1}{2025}$
\item $\lim_{n\to\infty}P(X_n=0)=\dfrac{1}{2024}$
\item $\lim_{n\to\infty}P(X_n=0)=\dfrac{1}{2023}$
\end{enumerate}\ans{2}

\textbf{Q.52} A mobile manufacturing company uses two brands of batteries for its mobiles. The life (in years) of batteries of Brand I follows an exponential distribution with the probability density function $f(x)=\begin{cases}e^{-x},&x>0\\0,&\text{otherwise}\end{cases}$, and that of Brand II follows a gamma distribution with the probability density function $g(x)=\begin{cases}\dfrac{x}{4}e^{-x/2},&x>0\\0,&\text{otherwise}\end{cases}$. The company uses the batteries of Brands I and II in proportion of 20\% and 80\% respectively. The probability that a randomly selected mobile has the battery life more than 2 years is
\begin{enumerate}[label=\arabic*.]
\item $\dfrac{13}{5}e^{-2}$
\item $\dfrac{1}{5}(e^{-2}+2e^{-1})$
\item $\dfrac{1}{5}(e^{-2}+8e^{-1})$
\item $\dfrac{1}{5}(4e^{-2}+2e^{-1})$
\end{enumerate}\ans{3}

\textbf{Q.53} Consider a discrete random variable $X$ with the probability mass function $P(X=0)=\dfrac{\theta}{3}$, $P(X=1)=1-\dfrac{\theta}{2}$, $P(X=2)=\dfrac{\theta}{6}$, where $\theta\in(0,1)$ is an unknown parameter. In a random sample of size 90 from this distribution, the observed counts for $X=0$, $1$ and $2$ are 20, 60 and 10, respectively. Then, the maximum likelihood estimate of $\theta$ is
\begin{enumerate}[label=\arabic*.]
\item $\dfrac{1}{3}$ \item $\dfrac{1}{2}$ \item $\dfrac{2}{3}$ \item $\dfrac{3}{4}$
\end{enumerate}\ans{3}

\textbf{Q.54} Let $X$ be a random sample of size 1 from the probability density function $f(x|\theta)=\begin{cases}\dfrac{3}{\theta^3}(\theta-x)^2,&0<x<\theta\\0,&\text{otherwise}\end{cases}$. If $\Bigl(\dfrac{X}{1-\lambda_1},\dfrac{X}{1-\lambda_2}\Bigr)$ is a confidence interval for $\theta$ with confidence coefficient $1-\alpha$, where $\lambda_i\in(0,1)$, $i=1,2$, $\lambda_1<\lambda_2$, and $\alpha\in(0,1)$, then which of the following is true?
\begin{enumerate}[label=\arabic*.]
\item $\lambda_2^2-\lambda_1^2=1-\alpha$
\item $\lambda_2^3-\lambda_1^3=1-\alpha$
\item $\lambda_2^2-\lambda_1^2=4(1-\alpha)$
\item $\lambda_2^3-\lambda_1^3=9(1-\alpha)$
\end{enumerate}\ans{2}

\textbf{Q.55} Let $X_1,X_2,\ldots,X_n$ be a random sample from a continuous distribution with the common probability density function $f(x|\theta)=\begin{cases}\dfrac{2^\theta\,\theta}{x^{\theta+1}},&x>2\\0,&\text{otherwise}\end{cases}$, where $\theta(>0)$ is an unknown parameter. Suppose $P(Y>\chi^2_{m,\beta})=\beta$, where $Y\sim\chi^2_m$. For testing $H_0:\theta=1$ against $H_1:\theta>1$, a uniformly most powerful test of size $\alpha$, $0<\alpha<1$, will reject $H_0$ if
\begin{enumerate}[label=\arabic*.]
\item $\sum_{i=1}^n\ln X_i>\tfrac{1}{2}\chi^2_{2n,\alpha}+n\ln2$
\item $\sum_{i=1}^n\ln X_i<\tfrac{1}{2}\chi^2_{2n,1-\alpha}+n\ln2$
\item $\sum_{i=1}^n\ln X_i>\chi^2_{n,\alpha}+n\ln2$
\item $\sum_{i=1}^n\ln X_i<\chi^2_{n,1-\alpha}+n\ln2$
\end{enumerate}\ans{1}

\textbf{Q.56} Suppose the distribution of $X$ given $\theta$ is normal with mean $\theta$ and variance 15. Further, let the prior (improper) distribution of $\theta$ be proportional to 1, $-\infty<\theta<\infty$. If the observed value of $X$ is 13, then which of the following is true?
\begin{enumerate}[label=\arabic*.]
\item Posterior mean $=$ Maximum likelihood estimate of $\theta$, Posterior variance $=\text{Var}(X|\theta)$
\item Posterior mean $=$ Maximum likelihood estimate of $\theta$, Posterior variance $<\text{Var}(X|\theta)$
\item Posterior mean $>$ Maximum likelihood estimate of $\theta$, Posterior variance $=\text{Var}(X|\theta)$
\item Posterior mean $>$ Maximum likelihood estimate of $\theta$, Posterior variance $<\text{Var}(X|\theta)$
\end{enumerate}\ans{1}

\textbf{Q.57} Consider the multiple linear regression model $y_i=\beta_0+\beta_1x_{1i}+\cdots+\beta_8x_{8i}+\epsilon_i$, $i=1,2,\ldots,29$, where $\epsilon_1,\epsilon_2,\ldots,\epsilon_{29}$ are independent and identically normally distributed with mean 0 and variance $\sigma^2$. Suppose the model is fitted using the method of least squares. If the calculated value of the $F$-statistic for testing the significance of regression is 2.50, then the possible values of $R^2$ and Adjusted $R^2$ are respectively
\begin{enumerate}[label=\arabic*.]
\item 0.30 and 0.10 \item 0.50 and 0.30 \item 0.50 and 0.16 \item 0.30 and $-0.10$
\end{enumerate}\ans{2}

\textbf{Q.58} Suppose $\begin{pmatrix}X_1\\X_2\\X_3\end{pmatrix}\sim N_3\!\!\left(\begin{pmatrix}0\\0\\0\end{pmatrix},\begin{pmatrix}1&0&0\\0&1&0\\0&0&1\end{pmatrix}\right)$ and $\begin{pmatrix}X\\Y\\Z\end{pmatrix}=\begin{pmatrix}3&0&0\\2&2&0\\4&0&4\end{pmatrix}\begin{pmatrix}X_1\\X_2\\X_3\end{pmatrix}$. Then the partial correlation coefficient $\rho_{YZ.X}$ is
\begin{enumerate}[label=\arabic*.]
\item $\dfrac{1}{2}$ \item $\dfrac{2}{3}$ \item $\dfrac{3}{4}$ \item $0$
\end{enumerate}\ans{4}

\textbf{Q.59} Suppose we want to estimate the population mean $\bar{Y}$ of a variable for a finite population of size 85, with 34 Statisticians and 51 Biologists. We consider the following sampling scheme: A stratified random sample with 2 strata of Statisticians (Stratum-1) and Biologists (Stratum-2), where 12 Statisticians and 15 Biologists are drawn from Stratum-1 and Stratum-2, respectively, using SRSWOR scheme. Denote $\bar{y}_S$, $\bar{y}_B$, and $\bar{y}$ as the mean of the variable among the Statistician sample, Biologist sample, and the combined sample, respectively. Which of the following is an unbiased estimator of $\bar{Y}$?
\begin{enumerate}[label=\arabic*.]
\item $\bar{y}$
\item $\dfrac{2\bar{y}_S+3\bar{y}_B}{5}$
\item $\dfrac{4\bar{y}_S+5\bar{y}_B}{9}$
\item $\dfrac{\dfrac{\bar{y}_S}{12}+\dfrac{\bar{y}_B}{15}}{\dfrac{1}{12}+\dfrac{1}{15}}$
\end{enumerate}\ans{2}

\textbf{Q.60} Solve the following linear programming problem: maximize $z=x+y$ subject to $5x+3y\le30$, $2x+6y\le25$, $2x-y\le8$, $x\ge0$, $y\ge0$. Then the optimal value of the objective function is
\begin{enumerate}[label=\arabic*.]
\item $\dfrac{45}{11}$ \item $\dfrac{74}{11}$ \item $\dfrac{85}{12}$ \item $\dfrac{25}{6}$
\end{enumerate}\ans{3}

%% ============================================================
\shead{PART C --- Mathematics MSQ (Q.61--Q.120)}
%% ============================================================
\begin{mdframed}[backgroundcolor=partbg,linewidth=0pt]
Marking: $+4.75$ fully correct (all correct options, no wrong), $-1.25$ for any wrong option selected. Attempt any 20 of 60.
\end{mdframed}
\vspace{0.5em}

\textbf{Q.61} Consider the sequences $(s_n)_{n\ge1}$ and $(t_n)_{n\ge1}$ defined by
\[s_n=\sum_{k=0}^n\dfrac{1}{(k!)^2}\quad\text{and}\quad t_n=\sum_{k=0}^n\binom{n}{k}\dfrac{(-1)^k}{n^k}.\]
Which of the following are true?
\begin{enumerate}[label=\arabic*.]
\item $\limsup_{n\to\infty}t_n\le\limsup_{n\to\infty}s_n$
\item $\limsup_{n\to\infty}t_n\le e$
\item $\liminf_{n\to\infty}s_n\ge e^2$
\item $\liminf_{n\to\infty}t_n\ge e$
\end{enumerate}\ans{1, 2}

\textbf{Q.62} What is the value of the limit $\displaystyle\lim_{n\to\infty}\dfrac{1}{n}\bigl[(n+1)(n+2)\cdots(n+n)\bigr]^{1/n}$?
\begin{enumerate}[label=\arabic*.]
\item $\dfrac{2}{e}$ \item $\dfrac{4}{e}$ \item $\log_e2-1$ \item $2\log_e2-1$
\end{enumerate}\ans{2}

\textbf{Q.63} Let $(f_n)_{n\ge1}$ be a sequence of real-valued functions on $\mathbb{R}$. Which of the following are true?
\begin{enumerate}[label=\arabic*.]
\item If each $f_n$ is uniformly continuous and $(f_n)_{n\ge1}$ converges to $f$ uniformly, then $f$ is uniformly continuous.
\item If each $f_n$ is bounded and $(f_n)_{n\ge1}$ converges to $f$ pointwise, then $f$ is bounded.
\item If each $f_n$ is bounded and continuous, $(f_n)_{n\ge1}$ converges pointwise to a bounded and continuous function $f$, then the convergence is uniform.
\item If each $f_n$ is differentiable and $(f_n)_{n\ge1}$ converges to $f$ uniformly, then $f$ is differentiable.
\end{enumerate}\ans{1}

\textbf{Q.64} Consider a dataset of $n$ observations given by $A=\{x_1,x_2,\ldots,x_n\}$. Create a new dataset, given by $B=\{x_1,x_2,\ldots,x_n,-x_1,-x_2,\ldots,-x_n\}$. Which of the following are always true?
\begin{enumerate}[label=\arabic*.]
\item Mean of the observations in dataset $B$ is 0
\item Median of the observations in dataset $B$ is 0
\item Variance of dataset $B\ge$ Variance of dataset $A$
\item Range of dataset $B\ge$ Range of dataset $A$
\end{enumerate}\ans{1, 2, 4}

\textbf{Q.65} Let $G$ be a group of order 8. Which of the following are necessarily true?
\begin{enumerate}[label=\arabic*.]
\item If there are no elements of order 4 in $G$, then $G$ is abelian.
\item If there are exactly two elements of order 4 in $G$, then $G$ is abelian.
\item If there are exactly six elements of order 4 in $G$, then $G$ is abelian.
\item There is an element of order 4 in $G$.
\end{enumerate}\ans{2}

\textbf{Q.66} Let $\alpha=\displaystyle\sum_{\ell=1}^{2025}\dfrac{a_\ell}{10^\ell}$ with $a_\ell\in\{0,1,\ldots,9\}$ for all $1\le\ell\le2025$ and $a_{2025}\neq0$. Let $(\beta_n)_{n\ge1}$ be a strictly increasing sequence of positive real numbers in $(0,1)$ that converges to $\alpha$. For each $n\ge1$, write $\beta_n=\displaystyle\sum_{\ell=1}^\infty\dfrac{b_{n,\ell}}{10^\ell}$ with $b_{n,\ell}\in\{0,1,\ldots,9\}$, for the infinite decimal expansion of $\beta_n$. Which of the following are necessarily true?
\begin{enumerate}[label=\arabic*.]
\item There exists a positive integer $N$ such that for all $n\ge N$, $b_{n,\ell}=a_\ell$ for all $1\le\ell\le2023$.
\item There exists a positive integer $N$ such that for all $n\ge N$, $b_{n,2025}=a_{2025}-1$.
\item $\beta_n$ is rational for infinitely many $n$.
\item There exists a positive integer $N$ such that for all $n\ge N$, $b_{n,2024}=a_{2024}-1$.
\end{enumerate}\ans{3, 4}

\textbf{Q.67} Let $\mathcal{F}$ be the set of all functions $f:[0,\infty)\to[0,\infty)$ that are Riemann integrable on $[0,t]$ for all $t\in[0,\infty)$. Which of the following are necessarily true?
\begin{enumerate}[label=\arabic*.]
\item If $f\in\mathcal{F}$ is a continuous function and $f(n)=0$ for all positive integers $n$, then $\lim_{x\to\infty}\int_0^xf(t)\,dt$ exists in $\mathbb{R}$.
\item If $f\in\mathcal{F}$ is uniformly continuous and $\lim_{x\to\infty}\int_0^xf(t)\,dt$ exists in $\mathbb{R}$, then $\lim_{x\to\infty}f(x)=0$.
\item There exists a function $f\in\mathcal{F}$ such that $\lim_{x\to\infty}\int_0^xf(t)\,dt$ exists in $\mathbb{R}$ but $\lim_{x\to\infty}f(x)\neq0$.
\item If $f\in\mathcal{F}$ is Lebesgue measurable and $\lim_{x\to\infty}\int_0^xf(t)\,dt=0$, then $f=0$ a.e.\ on $[0,\infty)$.
\end{enumerate}\ans{1, 3}

\textbf{Q.68} Consider the ODE $y''+r(x)y=0$, where $r(x)=\begin{cases}x^3\sin(1/x),&x\neq0\\0,&x=0\end{cases}$. Then which of the following are true?
\begin{enumerate}[label=\arabic*.]
\item There exists a non-trivial solution $\phi$ of ODE, and $\alpha,\beta\in\mathbb{R}$, $\alpha<\beta$, such that $\phi$ has infinitely many zeros in $[\alpha,\beta]$.
\item There exists a non-trivial solution $\phi$ of ODE, and a sequence $\{x_n\}$ such that $x_n\to\infty$ and $\phi(x_n)=0$ for all $n\in\mathbb{N}$.
\item If $\phi_1,\phi_2$ are two linearly independent solutions of ODE, then their Wronskian is a constant function on $\mathbb{R}$.
\item There exists a non-trivial solution of ODE which vanishes at most at one point in $(0,\infty)$.
\end{enumerate}\ans{2, 3}

\textbf{Q.69} Consider $X=\prod_{c\in(0,\infty)}[0,c]$ together with the product topology, where $[0,c]$ is equipped with the Euclidean topology. Which of the following are necessarily true?
\begin{enumerate}[label=\arabic*.]
\item $X$ is metrizable. \item $X$ is compact. \item $X$ is Hausdorff. \item $X$ is connected.
\end{enumerate}\ans{1, 3}

\textbf{Q.70} Let $X_1,X_2,\ldots,X_n$ ($n\ge2$) be a random sample from Uniform$\!\bigl[-\tfrac{3\theta}{2},\tfrac{\theta}{2}\bigr]$, $\theta>0$ unknown. Let $\bar{X}$, $X_{(1)}$ and $X_{(n)}$ denote the sample mean, smallest and largest order statistics, respectively. Which of the following are true?
\begin{enumerate}[label=\arabic*.]
\item The method of moments estimator of $\theta$ is $-2\bar{X}$
\item The method of moments estimator of $\theta$ is $2\bar{X}$
\item The MLE of $\theta$ is $\min\!\Bigl\{-\dfrac{2X_{(1)}}{3},2X_{(n)}\Bigr\}$
\item The MLE of $\theta$ is $\max\!\Bigl\{-\dfrac{2X_{(1)}}{3},2X_{(n)}\Bigr\}$
\end{enumerate}\ans{2, 3}

\textbf{Q.71} Let $\tau_E$ denote the Euclidean topology on $\mathbb{R}$. Which of the following are necessarily true?
\begin{enumerate}[label=\arabic*.]
\item $(\mathbb{R},\tau_E)$ is a normal space.
\item Let $\tau$ be a topology on $\mathbb{R}$. If the identity function is a continuous map from $(\mathbb{R},\tau_E)$ to $(\mathbb{R},\tau)$, then $(\mathbb{R},\tau)$ is a regular space.
\item Let $\tau$ be a topology on $\mathbb{R}$. If the identity function is a continuous map from $(\mathbb{R},\tau)$ to $(\mathbb{R},\tau_E)$, then $(\mathbb{R},\tau)$ is a Hausdorff space.
\item $\mathbb{R}$ is a regular space in the finite-complement topology.
\end{enumerate}\ans{3, 4}

\textbf{Q.72} Let $y(x)$ be the solution of the integral equation $y(x)=e^x+e+1+\tfrac{1}{3}\int_0^1y(t)\,dt$, and let $R\!\left(x,t,\tfrac{1}{3}\right)$ be the resolvent kernel associated to the IE. Then which of the following are true?
\begin{enumerate}[label=\arabic*.]
\item $R\!\left(0,1,\tfrac{1}{3}\right)=\dfrac{3}{2}$
\item $R\!\left(\tfrac{1}{2},\tfrac{1}{2},\tfrac{1}{3}\right)=\dfrac{2}{3}$
\item $y(1)=3e+1$
\item $y(1)=e+3$
\end{enumerate}\ans{1, 2, 3, 4}

\textbf{Q.73} Suppose $X|\theta\sim\text{Binomial}(7,\theta)$, $0<\theta<1$, and the prior distribution of $\theta$ is Beta$(\alpha,\beta)$ where $\alpha>0$ and $\beta>0$ are known. Which of the following \textbf{MAY NOT} be true?
\begin{enumerate}[label=\arabic*.]
\item Posterior mean of $\theta$ given $X=2$ is less than the prior mean
\item Posterior mean of $\theta$ given $X=3$ is greater than the prior mean
\item Posterior mean of $\theta$ given $X=4$ is not equal to the prior mean
\item Posterior mean of $\theta$ given $X=5$ is equal to the prior mean
\end{enumerate}\ans{1, 2}

\textbf{Q.74} Let $A\in M_4(\mathbb{C})$ be such that $A^2=I$ and $\text{Trace}(A)=0$. Which of the following are necessarily true?
\begin{enumerate}[label=\arabic*.]
\item The set $\{B\in M_4(\mathbb{C})\mid AB+BA=0\}$ is an 8-dimensional $\mathbb{C}$-vector space.
\item The set $\{B\in M_4(\mathbb{C})\mid AB-BA=0\}$ is an 8-dimensional $\mathbb{C}$-vector space.
\item There exists $B\in M_4(\mathbb{C})$ such that $AB+BA=0$, $B^2=I$ and $\text{Trace}(B)=0$.
\item If $B\in M_4(\mathbb{C})$ is such that $AB+BA=0$, then $B$ is diagonalisable.
\end{enumerate}\ans{1, 2, 4}

\textbf{Q.75} Consider an $M/M/2$ queuing system with birth rate $\lambda=4$ per minute, death rate $\mu=1$ per minute, and total capacity of 3 customers. Let $p_0$ and $p_3$ denote the long-run probabilities that the system will be empty and full, respectively. Which of the following are true?
\begin{enumerate}[label=\arabic*.]
\item $p_0+p_3=\dfrac{17}{29}$ \item $p_0>p_3$ \item $\dfrac{p_3}{p_0}=2$ \item $p_3-p_0=\dfrac{15}{29}$
\end{enumerate}\ans{1, 3}

\textbf{Q.76} For $n\ge3$, consider the space $M_n(\mathbb{C})$ of $n\times n$ complex matrices endowed with the inner product $\langle A,B\rangle=\text{Trace}(A^*B)$. For $0\le k\le n$, let $W_k=\{A\in M_n(\mathbb{C}):\langle A,B\rangle=0\text{ for all }B\in M_n(\mathbb{C})\text{ with rank }k\}$. Which of the following are necessarily true?
\begin{enumerate}[label=\arabic*.]
\item $W_0=\{0\}$ \item $W_1=\{0\}$ \item $W_2=\{0\}$ \item $W_n=\{0\}$
\end{enumerate}\ans{1, 3}

\textbf{Q.77} There are six persons. On each day of the year 2024 at least one of them borrowed money from another. A person $P$ is labeled \textit{troublesome} in calendar month $M$ of 2024 if $P$ borrowed from the same person at least twice during the month $M$. Which of the following are necessarily true?
\begin{enumerate}[label=\arabic*.]
\item In every calendar month of 2024, at least one person was troublesome.
\item In any two consecutive calendar months of 2024, at least one person was troublesome in at least one of the two calendar months.
\item At least one person was troublesome in two calendar months of 2024.
\item At least one person was troublesome in January 2024.
\end{enumerate}\ans{1, 2, 3, 4}

\textbf{Q.78} Let $V$ be a real vector space. Suppose $\{u,v,w,x,y\}\subseteq V$ is a spanning set of $V$ and that $\{u,v,x,y\}$ is linearly independent. Which of the following are necessarily true?
\begin{enumerate}[label=\arabic*.]
\item The dimension of $V$ is 4 or 5.
\item If $2u-3v+5w=0$, then $\{v,w,x,y\}$ is a basis of $V$.
\item If $u-6v+7w=0$, then the span of $\{v,w,x\}$ is a 2-dimensional vector space.
\item If $u+w=v+x$, then $\{u,v,w,y\}$ is a basis of $V$.
\end{enumerate}\ans{2, 4}

\textbf{Q.79} Suppose $y(x)$ is the extremal of the variational problem $J(y)=\int_0^1x^2(y')^2\,dx$ subject to $y(0)=0$, $y(1)=1$, $\int_0^1y^2\,dx=\tfrac{1}{7}$. Then which of the following are true?
\begin{enumerate}[label=\arabic*.]
\item $y'\!\left(\tfrac{1}{2}\right)=\dfrac{3}{4}$
\item $y'\!\left(\dfrac{1}{\sqrt{3}}\right)=1$
\item $y'\!\left(\tfrac{1}{3}\right)=\dfrac{1}{3}$
\item $y'\!\left(\tfrac{1}{4}\right)=\dfrac{1}{2}$
\end{enumerate}\ans{1, 3, 4}

\textbf{Q.80} Let $X_1,X_2,\ldots,X_{20}$ be a random sample from $N_{12}(\mu,\Sigma)$, where $\mu\in\mathbb{R}^{12}$ and $\Sigma$ is positive definite. Suppose $\bar{X}=\tfrac{1}{20}\sum_{i=1}^{20}X_i$ and $S=\tfrac{1}{19}\sum_{i=1}^{20}(X_i-\bar{X})(X_i-\bar{X})^T$. Then which of the following are true?
\begin{enumerate}[label=\arabic*.]
\item $19(\bar{X}^TS\bar{X})(\bar{X}^T\Sigma\bar{X})^{-1}\sim\chi^2_{19}$
\item $E(S^{-1})=\dfrac{19}{6}\Sigma^{-1}$
\item $S\sim W_{12}(19,19\Sigma)$
\item $\text{trace}(19\Sigma^{-1}S)\sim\chi^2_{228}$
\end{enumerate}\ans{2, 3}

\textbf{Q.81} Let $u(x,t)$ be the solution of the initial-boundary value problem $u_{tt}-u_{xx}+u_t=0$, $0<x<\pi$, $t>0$; $u(0,t)=0$, $u(\pi,t)=0$, $t>0$; $u(x,0)=\sin x$, $u_t(x,0)=0$, $0<x<\pi$. For $t\ge0$, define $E(t)=\int_0^\pi(u_t^2+u_x^2)\,dx$. Then which of the following are true?
\begin{enumerate}[label=\arabic*.]
\item $E(t)\ge\pi$ for all $t>0$. \item $E(t)\le\pi$ for all $t>0$.
\item $E(t)\ge\dfrac{\pi}{2}$ for all $t>0$. \item $E(t)\le\dfrac{\pi}{2}$ for all $t>0$.
\end{enumerate}\ans{4}

\textbf{Q.82} Let $S$ be the set of all $2\times2$ matrices $A$ such that the iterative sequence generated by the Gauss-Seidel method converges for every initial guess when employed to solve $A\begin{pmatrix}x_1\\x_2\end{pmatrix}=\begin{pmatrix}1\\2\end{pmatrix}$. Then which of the following are true?
\begin{enumerate}[label=\arabic*.]
\item $\begin{pmatrix}4&1\\2&3\end{pmatrix}\in S$
\item $\begin{pmatrix}1&2\\3&4\end{pmatrix}\in S$
\item $\begin{pmatrix}1&5\\1&10\end{pmatrix}\in S$
\item $\begin{pmatrix}-5&2\\1&-4\end{pmatrix}\in S$
\end{enumerate}\ans{2, 3, 4}

\textbf{Q.83} Consider a commutative ring $R$ with unity with at least five elements such that for any two elements $a,b\in R$, there exists $c\in R$ such that $a=bc$ or $b=ac$. Which of the following are necessarily true?
\begin{enumerate}[label=\arabic*.]
\item $R$ has a unique maximal ideal. \item Every finitely generated ideal of $R$ is principal.
\item $R$ is an integral domain. \item $R$ is a Euclidean domain.
\end{enumerate}\ans{1, 2}

\textbf{Q.84} Consider the following statements:\\
\textbf{(P)} The IVP $y'+y=e^{-x^2}$, $y(0)=0$, has a Taylor series solution about $x=0$.\\
\textbf{(Q)} The IVP $y'+y=r(x)$, $y(0)=0$, where $r(x)=\begin{cases}e^{-1/x^2},&x\neq0\\0,&x=0\end{cases}$, has a Taylor series solution about $x=0$.\\
Then which of the following are true?
\begin{enumerate}[label=\arabic*.]
\item (P) is true. \item (P) is FALSE.
\item (Q) is true. \item (Q) is FALSE.
\end{enumerate}\ans{4}

\textbf{Q.85} Let $F:\mathbb{R}^3\to\mathbb{R}$ be continuously differentiable such that $F(0,0,0)=-1$, $\dfrac{\partial F}{\partial x}(0,0,0)=2$, $\dfrac{\partial F}{\partial y}(0,0,0)=0$, and $\dfrac{\partial F}{\partial z}(0,0,0)=3$. For $\epsilon>0$, define $\Omega_\epsilon=(-\epsilon,\epsilon)\times(-\epsilon,\epsilon)\subset\mathbb{R}^2$. Which of the following are necessarily true?
\begin{enumerate}[label=\arabic*.]
\item There exist $\epsilon,\delta>0$ and continuously differentiable $g:\Omega_\epsilon\to(-\delta,\delta)$ such that $g(0,0)=0$ and $F(x,g(x,y),z)=-1$ for all $(x,y)\in\Omega_\epsilon$.
\item There exist $\epsilon,\delta>0$ and continuously differentiable $h:\Omega_\epsilon\to(-\delta,\delta)$ such that $h(0,0)=0$ and $F(x,h(x,z),z)=-1$ for all $(x,z)\in\Omega_\epsilon$.
\item There exist $\epsilon,\delta>0$ and continuously differentiable functions $k_1,k_2:(-\epsilon,\epsilon)\to(-\delta,\delta)$ such that $k_1(0)=k_2(0)=0$ and $F(x,k_1(x),k_2(x))=-1$ for all $x\in(-\epsilon,\epsilon)$.
\item There exist $\epsilon,\delta>0$ and continuously differentiable functions $j_1,j_2:(-\epsilon,\epsilon)\to(-\delta,\delta)$ such that $j_1(0)=j_2(0)=0$ and $F(j_1(z),j_2(z),z)=-1$ for all $z\in(-\epsilon,\epsilon)$.
\end{enumerate}\ans{1, 2}

\textbf{Q.86} For each positive integer $n$, let $f_n:[-1,1]\to\mathbb{R}$ be given by $f_n(x)=\begin{cases}-1,&-1\le x\le-\tfrac{1}{n}\\nx,&-\tfrac{1}{n}<x<\tfrac{1}{n}\\1,&\tfrac{1}{n}\le x\le1\end{cases}$. On which of the following intervals does the sequence $\{f_n\}_{n\ge1}$ converge uniformly?
\begin{enumerate}[label=\arabic*.]
\item $[0,1]$ \item $(0,1]$ \item $(10^{-2025},1)$ \item $(-10^{-2025},10^{-2025}]$
\end{enumerate}\ans{1, 2, 3}

\textbf{Q.87} Let $(X,d)$ be a metric space. For a non-empty subset $A$ of $X$, and $x\in X$, define $d(x,A)=\inf_{a\in A}d(x,a)$. Which of the following are necessarily true?
\begin{enumerate}[label=\arabic*.]
\item For all $x,y\in X$ and every non-empty subset $A$ of $X$, $d(x,A)-d(y,A)\le d(x,y)$.
\item For every non-empty subset $A$ of $X$, the function $x\mapsto d(x,A)$ is uniformly continuous on $X$.
\item A non-empty subset $A$ of $X$ is closed if and only if $d(x,A)>0$ for all $x\in X\setminus A$.
\item If $X$ is compact and $C_1,\ldots,C_k$ are non-empty closed sets of $X$ such that $\bigcap_{1\le i\le k}C_i=\emptyset$, then the minimum value of the function $x\mapsto\sum_{1\le i\le k}d(x,C_i)$ on $X$ is 0.
\end{enumerate}\ans{2}

\textbf{Q.88} Consider the following subset of real numbers: $A=\bigl\{(1+(-1)^n)n-\tfrac{1}{n}:n\text{ is a positive integer}\bigr\}$. Which of the following are true?
\begin{enumerate}[label=\arabic*.]
\item $A$ is bounded below but not bounded above.
\item $A$ is bounded above but not bounded below.
\item $\inf A=-1$ \item $\sup A=1$
\end{enumerate}\ans{2, 3, 4}

\textbf{Q.89} Let $f:(0,\infty)\to(0,\infty)$ be the function defined by $f(x)=\int_0^x\dfrac{\sqrt{t}}{1+t^2}\,dt$. Which of the following are true?
\begin{enumerate}[label=\arabic*.]
\item $f$ is a uniformly continuous function. \item $f$ is a bounded function.
\item There exists $x\in(0,\infty)$ such that $f(x)=0$. \item The derivative of $f$ is continuous.
\end{enumerate}\ans{1, 4}

\textbf{Q.90} Consider the following real-valued functions $F_1$ and $F_2$ defined on $\mathbb{R}$, given by $F_1(x)=\begin{cases}1,&x\ge1\\0,&\text{otherwise}\end{cases}$, $F_2(x)=\begin{cases}\int_0^xe^{-t}\,dt,&x\ge0\\0,&\text{otherwise}\end{cases}$. Define $F:\mathbb{R}\to[0,1]$ by $F(x)=\tfrac{2}{3}F_1(x)+\tfrac{1}{3}F_2(x)$ for all $x\in\mathbb{R}$. Which of the following are true?
\begin{enumerate}[label=\arabic*.]
\item $F$ is non-decreasing on $\mathbb{R}$ \item $\lim_{x\to\infty}F(x)=1$
\item $F$ is left-continuous on $\mathbb{R}$ \item $F$ is right-continuous on $\mathbb{R}$
\end{enumerate}\ans{1, 2}

\textbf{Q.91} Consider a $2^5$-factorial design with 8 blocks. (Block-treatment combination table as given in original question.) Which of the following are true?
\begin{enumerate}[label=\arabic*.]
\item $ABC$ is confounded with blocks \item $BCD$ is confounded with blocks
\item $CDE$ is confounded with blocks \item $ABDE$ is confounded with blocks
\end{enumerate}\ans{2, 3}

\textbf{Q.92} Let $\mathbb{N}$ denote the set of all positive integers. For $n\in\mathbb{N}$, let $f_n:[0,1]\to\mathbb{R}$ be given by $f_n(x)=\begin{cases}n(1-nx),&0\le x\le\tfrac{1}{n}\\0,&x>\tfrac{1}{n}\end{cases}$. Which of the following are necessarily true?
\begin{enumerate}[label=\arabic*.]
\item The sequence $\{f_n\}_{n\ge1}$ is uniformly bounded.
\item The sequence $\{f_n\}_{n\ge1}$ does not converge pointwise.
\item The sequence $\{f_n\}_{n\ge1}$ converges uniformly to the constant function 0.
\item The set $\{f_n:n\in\mathbb{N}\}$ is compact in $C[0,1]$ equipped with the supremum norm.
\end{enumerate}\ans{1, 2, 4}

\textbf{Q.93} Let $f$ be an entire function. Consider the function $g$ given by $g(z)=f(z)-\tfrac{1}{z}$ for $z\in\mathbb{C}\setminus\{0\}$. Which of the following are necessarily true?
\begin{enumerate}[label=\arabic*.]
\item The function $g$ has a pole at 0.
\item If $g(\alpha)=0$, then $|\alpha|\neq1$.
\item The function $g$ has only finitely many zeros.
\item $\max_{|z|=1}|g(z)|\ge1$
\end{enumerate}\ans{1, 3, 4}

\textbf{Q.94} Let $\ell^2$ denote the vector space of all square summable sequences $\{a_n\}_{n\ge1}$ of real numbers with the inner product $\langle\{a_n\},\{b_n\}\rangle=\sum_{n=1}^\infty a_nb_n$. Let $H=\bigl\{\{a_n\}\in\ell^2:|a_n|\le\tfrac{1}{n}\text{ for all positive integers }n\bigr\}$. Which of the following are true?
\begin{enumerate}[label=\arabic*.]
\item $H$ contains an orthonormal basis of $\ell^2$. \item $H$ is a linear subspace of $\ell^2$.
\item $H$ is a bounded subset of $\ell^2$. \item $H$ is a convex subset of $\ell^2$.
\end{enumerate}\ans{1}

\textbf{Q.95} Let $\{X_n:n\ge1\}$ be a sequence of i.i.d.\ random variables, where $X_1$ has an Exponential distribution with mean 1. Define $T_n=\max\{X_1,X_2,\ldots,X_n\}-\ln n$, $n\ge1$. Suppose $T_n\xrightarrow{d}Y$ as $n\to\infty$. Then which of the following are true?
\begin{enumerate}[label=\arabic*.]
\item $Y$ has a Double Exponential distribution with location parameter $\ln(\ln2)$ and scale parameter 1
\item Median of $Y=\ln(\ln2)$
\item $P(Y\le0)=e^{-1}$
\item The derivative of the CDF of $Y$ at $\ln3$ is $e^{-1/3}$
\end{enumerate}\ans{1, 4}

\textbf{Q.96} Let $f:[0,1]\to[0,1]$ be a monotonically increasing function. For any $\alpha\in(0,1)$, let $L_\alpha^+$ and $L_\alpha^-$ denote the right and left hand limits respectively. For $\alpha\in(0,1)$, if $L_\alpha^+$ and $L_\alpha^-$ exist, define $U_\alpha=\begin{cases}(L_\alpha^-,L_\alpha^+)&\text{if }L_\alpha^-<L_\alpha^+\\\emptyset&\text{if }L_\alpha^+\le L_\alpha^-\end{cases}$. Which of the following are true?
\begin{enumerate}[label=\arabic*.]
\item $L_\alpha^+$ and $L_\alpha^-$ exist for every $\alpha\in(0,1)$.
\item If $f$ is surjective, then $f$ is continuous.
\item If $f$ is Riemann integrable, then $f$ is continuous.
\item If the left and right hand limits exist at $\alpha,\beta\in(0,1)$, $\alpha\neq\beta$, then $U_\alpha\cap U_\beta=\emptyset$.
\end{enumerate}\ans{2, 3, 4}

\textbf{Q.97} Suppose that the transition probability matrix of a homogeneous Markov chain with state space $\{1,2,3,4\}$ is given by $P=\begin{pmatrix}\tfrac{1}{4}&\tfrac{3}{4}&0&0\\1&0&0&0\\\tfrac{1}{8}&0&\tfrac{7}{8}&0\\0&0&\tfrac{1}{9}&\tfrac{8}{9}\end{pmatrix}$. Which of the following are true?
\begin{enumerate}[label=\arabic*.]
\item State 2 is a positive recurrent state \item Mean recurrence time of state 1 is $\tfrac{7}{4}$
\item State 4 is a transient state \item State 3 is aperiodic and ergodic
\end{enumerate}\ans{1}

\textbf{Q.98} Let $X$ and $Y$ be two independent random variables such that $M_X(t)=e^{3(e^t-1)}$, $t\in\mathbb{R}$, and $M_Y(t)=\bigl(\tfrac{1}{2}e^{-3t}+\tfrac{1}{2}e^{3t}\bigr)^2$, $t\in\mathbb{R}$, respectively. Then which of the following are true?
\begin{enumerate}[label=\arabic*.]
\item $P(XY=0)=\dfrac{1+e^{-3}}{2}$ \item $E(X+Y)=3$
\item $\text{Var}(X+Y)=21$ \item $\text{Cov}(X+Y,X-Y)=0$
\end{enumerate}\ans{1, 2, 3}

\textbf{Q.99} Let $A$ be a non-zero $3\times3$ matrix with integer entries. Let $\lambda_i\in\mathbb{C}$, $1\le i\le3$, be all the eigenvalues of $A$. Which of the following are necessarily true?
\begin{enumerate}[label=\arabic*.]
\item There exists a cubic polynomial $f(X)\in\mathbb{Q}[X]$ such that $f(\lambda_i)=0$ for all $1\le i\le3$.
\item There exists a quadratic polynomial $f(X)\in\mathbb{Q}[X]$ such that $f(\lambda_i)=0$ for all $1\le i\le3$.
\item If $f(X)\in\mathbb{Q}[X]$ is such that $f(\lambda_i)=0$ for all $1\le i\le3$, then $f(A)=0$.
\item If $f(X)\in\mathbb{Q}[X]$ is a cubic polynomial such that $f(\lambda_i)=0$ for all $1\le i\le3$, then $f(A)=0$.
\end{enumerate}\ans{1, 2, 3, 4}

\textbf{Q.100} A mechanical system is described using generalized position $q$ and generalized momentum $p$. Let $Q$ and $P$ denote new generalized position and generalized momentum variables generated by the generating function $F(q,P)=q^2e^P$, and $Q$, $P$ are canonical coordinates. Let $G(p,Q)$ be a function such that $G(2,e)=0$, and it generates the same canonical coordinates $Q$, $P$. Then which of the following are true?
\begin{enumerate}[label=\arabic*.]
\item $G(p,Q)=-Q\!\left[1+\log_e\!\left(\dfrac{p^2}{4Q}\right)\right]$
\item $G(p,Q)=-pQ\!\left[1+\log_e\!\left(\dfrac{Q^2}{4p}\right)\right]$
\item $p=2qe^P$, $Q=q^2e^P$
\item $p=-2qe^P$, $Q=-q^2e^P$
\end{enumerate}\ans{1, 2, 3}

\textbf{Q.101} Let $\mathbb{D}=\{z\in\mathbb{C}:|z|<1\}$ and $f:\mathbb{D}\to\mathbb{D}$ be a holomorphic function which satisfies $f(-1/2)=0$. Which of the following are necessarily true?
\begin{enumerate}[label=\arabic*.]
\item $|f(-1/5)|\le1/5$ \item $|f(-1/5)|\le1/3$
\item $|f'(-1/2)|\le1/2$ \item $|f'(-1/2)|\le4/3$
\end{enumerate}\ans{1, 2, 4}

\textbf{Q.102} For $t\in[0,2\pi]$, let $\gamma_1(t)=e^{it}$ and $\gamma_2(t)=1+i+e^{it}$. Which of the following are true?
\begin{enumerate}[label=\arabic*.]
\item $\displaystyle\int_{\gamma_1}\dfrac{\mathrm{d}z}{z\sin z}=0$
\item $\displaystyle\int_{\gamma_1}\dfrac{\mathrm{d}z}{z\sin z}=2\pi i$
\item $\displaystyle\int_{\gamma_2}\dfrac{\mathrm{d}z}{z\sin z}=2\pi i$
\item $\displaystyle\int_{\gamma_2}\dfrac{\mathrm{d}z}{z\sin z}=0$
\end{enumerate}\ans{1, 2, 3, 4}

\textbf{Q.103} Let $V$ be a finite-dimensional $\mathbb{R}$-vector space and $T:V\to V$ a linear operator such that $T^2$ is diagonalizable over $\mathbb{R}$. Which of the following are necessarily true?
\begin{enumerate}[label=\arabic*.]
\item If $T$ is not diagonalizable over $\mathbb{R}$, then $T^2$ has an eigenvalue $\le0$.
\item If $T^2$ has only negative eigenvalues, then $\dim V$ is an even integer.
\item If $T^2$ has only non-negative eigenvalues, then $T$ is diagonalizable over $\mathbb{R}$.
\item For each non-zero $v\in V$, $\{v,Tv,T^2v\}$ is linearly dependent.
\end{enumerate}\ans{2, 3, 4}

\textbf{Q.104} Let $X_1,X_2,\ldots,X_7$ be a random sample drawn from a continuous distribution with unknown unique median $M$. The null hypothesis $H_0:M=2$ is tested against $H_1:M>2$ at significance level 0.05 using the right-tailed Sign test statistic $K$ (number of observations greater than 2). If the observed sample is $-3,-6,1,9,4,10,12$, which of the following are true?
\begin{enumerate}[label=\arabic*.]
\item $H_0$ is rejected
\item The $p$-value of the test is greater than 0.01
\item Under $H_0$, $E(K)=3.5$
\item If $X_{(i)}$ denotes the $i$-th smallest observation of the sample, then $[X_{(2)},X_{(6)}]$ is a confidence interval for $M$ with confidence coefficient at least 0.95
\end{enumerate}\ans{1, 3}

\textbf{Q.105} Suppose $y_1(x)$ and $y_2(x)$ are two linearly independent solutions of the differential equation $x^2y''+\dfrac{(1+\sin x)x}{2}y'-\dfrac{3}{2}(\cos x)y=0$, $x>0$, satisfying $y_2(0)=0$. Then which of the following are true?
\begin{enumerate}[label=\arabic*.]
\item $\lim_{x\to0^+}\dfrac{y_2(x)}{x^2y_1(x)}$ exists.
\item $\lim_{x\to0^+}\dfrac{y_2(x)}{xy_1(x)}$ exists.
\item $\lim_{x\to0^+}\dfrac{xy_1(x)}{y_2(x)}$ does NOT exist.
\item $\lim_{x\to0^+}\dfrac{x^2y_1(x)}{y_2(x)}$ exists.
\end{enumerate}\ans{1, 2, 3, 4}

\textbf{Q.106} Suppose two fair dice are thrown independently at random. Let $X$ and $Y$ be the numbers on the upper face of the first die and that of the second die, respectively. Then which of the following are true?
\begin{enumerate}[label=\arabic*.]
\item $P(X-Y=0\mid X+Y=2)=P(X-Y=0\mid X+Y=12)$
\item $E\!\left(\dfrac{X-Y}{X+Y}\right)=0$
\item $\text{Cov}(X+Y,X-Y)=0$
\item $(X+Y)$ and $(X-Y)$ are independent
\end{enumerate}\ans{1, 3}

\textbf{Q.107} Let $V$ be a 4-dimensional complex vector space and $A$ a linear operator on $V$. Which of the following are necessarily true?
\begin{enumerate}[label=\arabic*.]
\item There exist $\lambda\in\mathbb{C}$ and a non-zero $v\in V$ such that $Av=\lambda v$.
\item There exist $\lambda,\mu\in\mathbb{C}$ and linearly independent vectors $v,w\in V$ such that $Av=\lambda v$ and $Aw=\mu w$.
\item There exist $\lambda,\mu,\delta\in\mathbb{C}$ and linearly independent vectors $v,w\in V$ such that $Av=\lambda v$ and $Aw=\mu v+\delta w$.
\item There exists a three-dimensional subspace $W$ such that $Aw\in W$ for all $w\in W$.
\end{enumerate}\ans{1, 2}

\textbf{Q.108} Let $p\ge3$ be a prime number and $f(x)\in\mathbb{Q}[x]$ an irreducible polynomial of degree $p$. Suppose that $a_1,\ldots,a_p\in\mathbb{C}$ are the roots of $f$ and that $a_1\notin\mathbb{R}$, $a_2\notin\mathbb{R}$, and $a_i\in\mathbb{R}$ for all $3\le i\le p$. Let $K=\mathbb{Q}(a_1,\ldots,a_p)$ be the subfield of $\mathbb{C}$ generated by the roots of $f$. Consider the Galois group $G$ of $K$ over $\mathbb{Q}$ as a subgroup of $S_p$, the group of permutations of $\{a_1,\ldots,a_p\}$. Which of the following are true?
\begin{enumerate}[label=\arabic*.]
\item The transposition $(1\,2)$ belongs to $G$. \item $|G|$ is divisible by $p$.
\item A $p$-cycle belongs to $G$. \item $G=S_p$
\end{enumerate}\ans{1, 2, 4}

\textbf{Q.109} Let $V$ be a real vector space and $L(V)$ denote the space of linear operators on $V$. Let $T\in L(V)$ be a non-zero operator such that $T^2=T$. Consider the subspace $W$ of $L(V)$ spanned by $\{I,T^n:n\text{ is a positive integer}\}$. Which of the following are necessarily true?
\begin{enumerate}[label=\arabic*.]
\item The set $\{S\in W:S^2=S\}$ contains exactly 2 elements.
\item $\dim_\mathbb{R}(W)=2$
\item If $U\in L(V)$ is such that $U^2=U$ and $(T+U)^2=T+U$, then $TU=0$.
\item If $U\in L(V)$ is such that $U^2=U$ and $(T-U)^2=T-U$, then $(TU)^2=TU$.
\end{enumerate}\ans{3, 4}

\textbf{Q.110} Let $G$ be a finite non-abelian group. Which of the following are necessarily true?
\begin{enumerate}[label=\arabic*.]
\item If $d$ is a positive integer that divides $|G|$, then $G$ has a subgroup of order $d$.
\item The map $f:G\times G\to G$ given by $f(a,b)=ab$ is not a group homomorphism.
\item Suppose that for every positive integer $d$ that divides $|G|$, there exists a subgroup of $G$ of order $d$. Then $G$ has at least three normal subgroups.
\item $|G|\neq16$
\end{enumerate}\ans{1, 3, 4}

\textbf{Q.111} Consider a series system comprising four components, having independent and identically distributed lifetimes with hazard rate $\lambda(t)=\dfrac{1}{1+t}$, $t>0$, and survival function $S(t)$, $t>0$. If $Y$ denotes the lifetime of the series system, then which of the following are true?
\begin{enumerate}[label=\arabic*.]
\item Cumulative hazard function of each component is $H(t)=2\ln(1+t)$, $t>0$
\item $S(t)=\dfrac{1}{(1+t)^2}$, $t>0$
\item $P(Y<\tfrac{1}{2})=\bigl\{S(\tfrac{1}{2})\bigr\}^4$
\item $P(Y<\tfrac{1}{2})=\dfrac{65}{81}$
\end{enumerate}\ans{1, 3}

\textbf{Q.112} Let $a,b$ be distinct positive real numbers. Consider the function $f:\mathbb{R}^2\to\mathbb{R}$ given by $f(x,y)=\begin{cases}\dfrac{(ax+by)^2}{ax^2+by^2}&\text{if }(x,y)\neq(0,0)\\0&\text{if }(x,y)=(0,0)\end{cases}$. Which of the following are necessarily true?
\begin{enumerate}[label=\arabic*.]
\item $\lim_{(x,y)\to(0,0)}f(x,y)$ does not exist.
\item The partial derivatives of $f$ at $(0,0)$ do not exist.
\item $\lim_{x\to0}f(x,0)=\lim_{y\to0}f(0,y)$
\item $f$ is differentiable at $(0,0)$.
\end{enumerate}\ans{1, 2, 3}

\textbf{Q.113} Let $X_1,X_2,\ldots,X_n$ ($n>5$) be independent random variables such that $X_t=\alpha+\alpha^2t+\epsilon_t$, $t=1,\ldots,n$, where $\epsilon_1,\epsilon_2,\ldots,\epsilon_n$ are i.i.d.\ $N(0,\sigma^2)$ random variables, $\alpha\in\mathbb{R}$ and $\sigma>0$ are unknown parameters. Which of the following are true?
\begin{enumerate}[label=\arabic*.]
\item $(X_1,X_2,\ldots,X_n)$ is a sufficient statistic for $(\alpha,\sigma)$
\item $\Bigl(\sum_{t=1}^nX_t,\sum_{t=1}^ntX_t,\sum_{t=1}^nt^2X_t^2\Bigr)$ is a jointly minimal sufficient statistic for $(\alpha,\sigma)$
\item $X_4-X_3-X_2+X_1$ is an ancillary statistic
\item $\dfrac{X_4+X_1-X_2-X_3}{X_5+X_2-X_3-X_4}$ is an ancillary statistic
\end{enumerate}\ans{4}

\textbf{Q.114} Suppose $y(x)$ is the extremal of the variational problem $J(y)=\int_0^1\bigl[(y')^2\sin x+(2\cos x)y\bigr]\,dx$ subject to $y(0)=0$, $y(1)=1$. Then which of the following are true?
\begin{enumerate}[label=\arabic*.]
\item $y\!\left(\tfrac{1}{2}\right)=1$ \item $y'(0)=1$
\item $y\!\left(\tfrac{1}{4}\right)=2$ \item $y'\!\left(\tfrac{1}{2}\right)=1$
\end{enumerate}\ans{1, 4}

\textbf{Q.115} Let $Y=X\beta+\epsilon$ be a multiple linear regression model with $p$ regressors and an intercept, where $\beta=(\beta_0,\beta_1,\ldots,\beta_p)^T$ and the random error $\epsilon\sim N_n(0,\sigma^2I)$, $\sigma>0$ and $n>p+1$. Let $\hat{\beta}$ be the least squares estimator of $\beta$. Let the total sum of squares (corrected), sum of squares due to regression/model and sum of squares due to error/residual, based on $\hat{\beta}$, be denoted by $Y^TAY$, $Y^TBY$ and $Y^TCY$ respectively, so that $Y^TAY=Y^TBY+Y^TCY$. Which of the following are always true?
\begin{enumerate}[label=\arabic*.]
\item $\dfrac{Y^TAY}{\sigma^2}$ follows a central $\chi^2$ distribution with $(n-1)$ degrees of freedom.
\item $\dfrac{Y^TBY}{\sigma^2}$ follows a central $\chi^2$ distribution with $p$ degrees of freedom if $\beta_1=\cdots=\beta_p=0$.
\item $\dfrac{Y^TBY}{Y^TCY}$ follows a central $F$ distribution with $(p,n-p)$ degrees of freedom.
\item $Y^TBY$ and $Y^TCY$ are independently distributed if and only if $\beta_1=\cdots=\beta_p=0$.
\end{enumerate}\ans{1, 4}

\textbf{Q.116} Let $\lambda\in\mathbb{R}$ be such that the integral equation $y(x)=\lambda\int_{-1}^1(5xt^3+4x^2t+3xt)y(t)\,dt$ admits a non-trivial solution $y(x)$ such that $y(1)=\tfrac{5}{2}$. Then which of the following are true?
\begin{enumerate}[label=\arabic*.]
\item $y(0)+y'(0)=\dfrac{3}{2}$
\item $y\!\left(\tfrac{1}{2}\right)+y'\!\left(\tfrac{1}{2}\right)=\dfrac{7}{2}$
\item $y(-1)+y'(-1)=-1$
\item $y\!\left(\tfrac{1}{3}\right)+y'\!\left(\tfrac{1}{3}\right)=\dfrac{14}{9}$
\end{enumerate}\ans{1, 3}

\textbf{Q.117} Let $\{x_n\}$ be a convergent iterative sequence generated by Newton-Raphson method for solving $\sin x-1=0$ such that $x_n\to\dfrac{\pi}{2}$ as $n\to\infty$. For $n\in\mathbb{N}$, let $e_n=x_n-\dfrac{\pi}{2}$. Let $p>0$ be such that $\lim_{n\to\infty}\dfrac{|e_{n+1}|}{|e_n|^p}$ exists and is non-zero. Then which of the following are true?
\begin{enumerate}[label=\arabic*.]
\item $p=1$ \item $\lim_{n\to\infty}\dfrac{|e_{n+1}|}{|e_n|^p}=\dfrac{1}{2}$
\item $p=2$ \item $\lim_{n\to\infty}\dfrac{|e_{n+1}|}{|e_n|^p}=1$
\end{enumerate}\ans{1}

\textbf{Q.118} Consider the boundary value problem (BVP) $u_{xx}+u_{yy}=0$ in $\Omega=\{(x,y)\in\mathbb{R}^2:x^2+y^2<1\}$, $u(x,y)=e^{x+y}$ on $\partial\Omega=\{(x,y)\in\mathbb{R}^2:x^2+y^2=1\}$. Then which of the following are true?
\begin{enumerate}[label=\arabic*.]
\item There exists a unique solution to BVP.
\item The BVP does NOT have a solution.
\item There exists a solution $u$ to BVP such that $u(x,y)=\dfrac{1+e}{2}$ for some $(x,y)\in\Omega\cup\partial\Omega$.
\item There exists a solution $u$ to BVP such that $u(x,y)=\dfrac{1+e^3}{2}$ for some $(x,y)\in\Omega\cup\partial\Omega$.
\end{enumerate}\ans{1, 2, 3}

\textbf{Q.119} Consider the simple linear regression model $Y_i=\beta x_i+\epsilon_i$, $i=1,2,\ldots,n$, where $\beta>0$, $\sum_{i=1}^nx_i^2>0$, and the uncorrelated errors $\epsilon_i$ have zero mean and finite variance $\sigma^2(>0)$. Let $\tilde{\beta}_1=\sum_{i=1}^na_i^*Y_i$ minimise $E\!\bigl(\sum_{i=1}^na_iY_i-\beta\bigr)^2$ with respect to scalars $a_1,\ldots,a_n$. Let $\tilde{\beta}_2$ be the ordinary least squares estimator of $\beta$. Which of the following are true?
\begin{enumerate}[label=\arabic*.]
\item $\text{Var}(\tilde{\beta}_1)>\text{Var}(\tilde{\beta}_2)$, $E(\tilde{\beta}_1-\beta)^2>E(\tilde{\beta}_2-\beta)^2$
\item $\text{Var}(\tilde{\beta}_1)<\text{Var}(\tilde{\beta}_2)$, $E(\tilde{\beta}_1-\beta)^2<E(\tilde{\beta}_2-\beta)^2$
\item $\text{Var}(\tilde{\beta}_1)>\text{Var}(\tilde{\beta}_2)$, $E(\tilde{\beta}_1)<E(\tilde{\beta}_2)$
\item $\text{Var}(\tilde{\beta}_1)<\text{Var}(\tilde{\beta}_2)$, $E(\tilde{\beta}_1)>E(\tilde{\beta}_2)$
\end{enumerate}\ans{1, 3}

\textbf{Q.120} Let $f:\mathbb{C}\to\mathbb{C}$ be defined by $f(z)=\sin^2z+\cos^2|z|$. Which of the following are true?
\begin{enumerate}[label=\arabic*.]
\item $f$ is a real valued function.
\item $f(z)=1$ for all $z\in\mathbb{C}$.
\item $f$ is not an entire function.
\item $f$ has finitely many zeros on the imaginary axis.
\end{enumerate}\ans{2, 4}

\vspace{1em}\hrule
\begin{center}
\textbf{END OF CSIR-UGC NET JUNE 2025 --- MATHEMATICAL SCIENCES}
\end{center}

\end{document}
